\begin{thebibliography}{19}
\small

\bibitem{nhis} National Center for Health Statistics. National Health Interview Survey, Sample Adult files 2019--2023. Centers for Disease Control and Prevention. \url{https://www.cdc.gov/nchs/nhis/}

\bibitem{epa} US Environmental Protection Agency. AirData: annual concentration by monitor, 2019--2023. \url{https://aqs.epa.gov/aqsweb/airdata/download_files.html}

\bibitem{brfss} CDC Behavioral Risk Factor Surveillance System. BRFSS prevalence data, current smokers, 2019--2023. \url{https://data.cdc.gov/resource/dttw-5yxu}

\bibitem{uscs} CDC and National Cancer Institute. US Cancer Statistics, incidence by area, 1999--2023. \url{https://www.cdc.gov/cancer/uscs/}

\bibitem{kaggle} Kaggle user mysarahmadbhat. Lung Cancer (survey) dataset, 309 respondents; compared through the identical public copy at github.com/Akash6181/LungCancerSurvey. The rejected file \texttt{lung\_cancer\_combined\_22457.csv} was provided at the start of the project; its original provenance is not documented. \url{https://www.kaggle.com/datasets/mysarahmadbhat/lung-cancer}

\bibitem{uspstf} US Preventive Services Task Force (2021). Screening for lung cancer: recommendation statement. \emph{JAMA}, 325(10), 962--970. doi:10.1001/jama.2021.1117

\bibitem{tammemagi} Tammem\"agi, M.\,C. et al. (2013). Selection criteria for lung-cancer screening. \emph{New England Journal of Medicine}, 368(8), 728--736. doi:10.1056/NEJMoa1211776

\bibitem{holm} Holm, S. (1979). A simple sequentially rejective multiple test procedure. \emph{Scandinavian Journal of Statistics}, 6(2), 65--70.

\bibitem{cohen} Cohen, J. (1988). \emph{Statistical Power Analysis for the Behavioral Sciences} (2nd ed.). Lawrence Erlbaum.

\bibitem{crispdm} Chapman, P. et al. (2000). \emph{CRISP-DM 1.0: Step-by-step data mining guide}. SPSS Inc.

\bibitem{kapoor} Kapoor, S. and Narayanan, A. (2023). Leakage and the reproducibility crisis in machine-learning-based science. \emph{Patterns}, 4(9), 100804. doi:10.1016/j.patter.2023.100804

\bibitem{tripod} Collins, G.\,S. et al. (2024). TRIPOD+AI statement: updated guidance for reporting clinical prediction models that use regression or machine learning methods. \emph{BMJ}, 385, e078378. doi:10.1136/bmj-2023-078378

\bibitem{scipy} Virtanen, P. et al. (2020). SciPy 1.0: fundamental algorithms for scientific computing in Python. \emph{Nature Methods}, 17, 261--272.

\bibitem{sklearn} Pedregosa, F. et al. (2011). Scikit-learn: machine learning in Python. \emph{Journal of Machine Learning Research}, 12, 2825--2830.

\bibitem{statsmodels} Seabold, S. and Perktold, J. (2010). Statsmodels: econometric and statistical modeling with Python. \emph{Proceedings of the 9th Python in Science Conference}, 92--96.

\bibitem{mantel} Mantel, N. and Haenszel, W. (1959). Statistical aspects of the analysis of data from retrospective studies of disease. \emph{Journal of the National Cancer Institute}, 22(4), 719--748.

\bibitem{larose} Larose, D.\,T. and Larose, C.\,D. (2015). \emph{Data Mining and Predictive Analytics} (2nd ed.). Wiley.

\bibitem{jolliffe} Jolliffe, I.\,T. (2002). \emph{Principal Component Analysis} (2nd ed.). Springer.

\bibitem{louati} Louati, D. (2026). CS434 Data Science, lecture slides, Chapters 1 (Foundations, CRISP-DM), 2 (Data gathering and ETL), 4 (Exploratory data analysis), 5 (Feature engineering) and 6 (Descriptive modelling); and Data Science Midterm Project -- Guidelines and Evaluation Rubric. MedTech -- Mediterranean Institute of Technology.

\end{thebibliography}
